% Save the original definition
\let\oldtextbf\textbf

% Redefine \textbf to include color
\renewcommand{\textbf}[1]{\oldtextbf{\textcolor{RoyalBlue}{#1}}}

\title{Grundlagenfach Informatik, Schuljahr 2025/2026}
\author{Cyril Wendl}

\begin{document}

\lhead[ \leftmark ]{Informatik, Semesterplanung}
\rhead[Informatik]{Cyril Wendl, \today, Winterthur}

\maketitle

\section*{Thematische Übersicht}
Willkommen in der vielfältigen Welt der Informatik! \textbf{Mein Hauptziel}: Ihnen \textbf{Freude an der Informatik} vermitteln! 

Im \textbf{Grundlagenfach Informatik} werden Sie viele Aspekte dieser spannenden Wissenschaft kennenlernen:

\begin{itemize}
	\item \textbf{Kodierungen}: Wie funktionieren Computer und wie sind diese aufgebaut?
	\item \textbf{Programmieren}: Wie können wir dem Computer spannende und komplexe Aufgaben geben, welche unseren Alltag erleichtern und unsere Kreativität beflügeln?
	\item \textbf{Verschlüsselung}: Wie können wir unsere Daten schützen, so dass keine unbefugten Personen darauf Zugriff haben?
	\item \textbf{Kompression}: Wie können wir die riesigen Datenberge, die jeden Tag wachsen, effizient speichern?
	\item \textbf{Datenintegrität}: Wie können wir sicherstellen, dass es keine Fehler bei der Übertragung von Daten gibt?
	\item \textbf{Datenbanken}: Wie können wir effizient auf Daten zugreifen, diese verändern und zusammenfügen, um spannende neue Einsichten zu gewinnen?
	\item \textbf{Netzwerke}: Wie funktioniert das Internet?
	\item \textbf{Robotik}: Wie können wir Roboter programmieren, welche selbstständig komplexe Aufgaben erledigen?
	\item \textbf{Game-Programmierung}: Wie können wir spannende Spiele programmieren, die uns unterhalten und unser Verständnis für die digitale Welt fördern?
	\item \textbf{Machine Learning und Klassifikation}: Wie können wir dem Computer beibringen, automatische Vorhersagen für neue Daten zu treffen (Wetter, Pflanzenart, Krankheiten, Börsenkurse etc.)?
\end{itemize}

Einige dieser Themen werden wir aus Zeitgründen nur anschneiden. Falls Sie sich weiter in diese Themen vertiefen möchten, empfehlen wir Ihnen das Ergänzungsfach Informatik, in welchem Sie viele dieser Themen vertiefen und mit spannenden Projekten umsetzen können.

Nicht zuletzt werden wir während des gesamten Unterrichts immer wieder über \textbf{gesellschaftliche Folgen} der Digitalisierung sprechen und neuste Trends kritisch gemeinsam betrachten.

Falls es \textbf{weitere Themen} gibt, welche Sie interessieren, dürfen Sie diese sehr gerne mit mir diskutieren, so dass ich diese im Unterricht einbringen kann.

\section*{Grundwerte}
Ich wünsche mir eine \textbf{entspannte} und \textbf{freundliche} Klassen-Atmosphäre, in der sich alle wohl fühlen. Dazu gehört für mich eine Prise (Selbst-)Ironie und \textbf{Humor}, allerdings auch \textbf{Ernsthaftigkeit}, \textbf{Neugier} und \textbf{Offenheit}, um sich auf Neues einzulassen.

\textbf{Pünktlichkeit und Verlässlichkeit} sind mir sehr wichtig: Damit Sie möglichst effektiv lernen können, ist es wichtig, dass Sie die wenigen erteilten Hausaufgaben pflichtbewusst erledigt und Termine notieren. Auch ich werde mich an diesem Massstab messen und werde versuchen, Ihre Prüfungen stets zeitnah und gewissenhaft zu korrigieren und bewerten.

Ich schätze Ihr \textbf{Feedback} auf meinen Unterricht sehr, da mir dies hilft, mich stetig zu verbessern. Daher werde ich Sie regelmässig um ihr freiwilliges, anonymes Feedback zum Unterricht bitten. Sie können mir stets auch persönliche Rückmeldungen geben. Ich freue mich über jegliches konstruktives und möglichst konkret umsetzbares Feedback von Ihnen! Feedback, welches Sie mir geben, dient Ihnen nicht zuletzt auch zur Aufbesserung Ihrer Semesternote (siehe Benotungsrichtlinien).

Allgemein setze ich alles daran, den Unterricht für Sie möglichst spannend und kurzweilig zu machen -- ich kann jedoch nicht garantieren, dass dies immer gelingt. Im Gegensatz erwarte ich von Ihnen ebenfalls eine interessierte und aktive Teilnahme am Unterricht, sowie die Bereitschaft, sich auf Neues einzulassen.

\begin{quote}
\textit{Achte auf Deine Gedanken, denn es liegt in der Natur des Gedankes, Wirklichkeit zu werden.}
\end{quote}

\section*{Regeln}

Nebst offensichtlichen Grundregeln der Schule gelten für den Informatikunterricht folgende Regeln:

\begin{itemize}
	\item[\faLaptop] Bitte bringen Sie in jede Unterrichtsstunde ihren \textbf{voll aufgeladenen Laptop} mit, sowie ein \textbf{Ladekabel} dafür.
	\item[\faTint] Ich bitte Sie freundlich, während dem Unterricht \textbf{nur Wasser} zu trinken und Essen sowie andere Getränke auf die Pause zu verschieben.
	\item[\faMobile] Bitte \textbf{keine Handys} im Unterricht, da diese den Unterricht stören und Sie ablenken.
	\item[\faHeadphones] Bitte \textbf{keine Kopfhörer, Hüte oder Sonnenbrillen} im Unterricht, da ich gerne mit Ihnen direkt kommunizieren möchte.
	\item[\faRestroom] \textbf{WC} sowie Wasser auffüllen nur in begründeten Ausnahmefällen.
	\item[\faClock] Ich bitte Sie, jeweils \textbf{pünktlich} im Unterricht zu erscheinen.
\end{itemize}

\begin{notizTemplate}[title={Weitere, gemeinsam vereinbarte Grundwerte und Regeln}, height=5cm]
	
\end{notizTemplate}

\begin{notizTemplate}[title={\emoji{pen} Unterschriften},height=7cm,colframe=RoyalBlue,colbacktitle=RoyalBlue]

\end{notizTemplate}



\end{document}